\section{Constructing the infinite innovation equation}
\label{sec:innovations}

The posterior fields describe how the observations affect the spins.
We now turn that description into a stochastic equation driven by
independent Brownian motions.  This is the equation that will allow us
to compare two observation processes driven by the same noise.
Throughout the section,
\[
 b=d-1,\qquad 0<\theta=\tanh\beta\le b^{-1/2}.
\]
We use the raw observation versions from Lemma~\ref{lem:posterior}.
In particular, at this stage a posterior identity holds almost surely
at each fixed time, and therefore outside a
$\dd t\,\dd\P$-null set; continuity of the infinite filter has not yet
been established.

For each vertex define its \emph{innovation} by
\begin{equation}
 \label{eq:innovations-definition}
 W_v(t)=X_v(t)-\int_0^t m_v(s)\dd s.
\end{equation}
The integrand is a bounded, progressively measurable observation
function.  Thus this formula defines continuous paths for every
observation input.  Their values through time $t$ use only the
observation history through $t$.  The time integrals are Borel
functions of that history, so $X\mapsto W$ is a causal Borel map on
the countable product of continuous path spaces.  The centered-ball
definitions of the posterior means also show that this map commutes
with graph relabeling.

The innovation construction is classical in nonlinear filtering
\cite{FujisakiKallianpurKunita1972} and is the starting point of the
Brownian Ising construction of Nam, Sly, and Zhang~\cite{NSZ}.  Here
we must justify its equation for countably many observations.  We will
first prove that $W$ is Brownian, then identify the entire posterior
martingale using an exact energy limit.  Neither step assumes that
$W$ determines $X$.

\begin{proposition}[Infinite innovation equations]
 \label{prop:innovations}
In the planted observation law $X_v(t)=t\sigma_v+B_v(t)$, every finite
subvector of $(W_v)_{v\in\mathcal V}$ is a standard vector Brownian
motion in both the raw observation filtration $(\mathcal F_t^0)$ and
its usual augmentation $(\mathcal F_t)$.  In particular $W$ has product
Wiener law.

The posterior means have continuous versions, starting at zero, and
\begin{equation}
 \label{eq:magnetization-sde}
 m_v(t)=\sum_{u\in\mathcal V}\int_0^t C_{vu}(s)\dd W_u(s).
\end{equation}
At each fixed time these versions are
$\E[\sigma_v\mid\mathcal F_t]$, and $C_{vu}$ is the covariance
conditional on $\mathcal F_t$.  The site fields and the component
fields can be chosen as jointly Borel, progressively measurable
functions of time and observation history that are continuous in time
almost surely.  For every input, the chosen function
$h_{v\to w}(t)$ uses only observations in $T_{v\to w}$ through time
$t$.  Almost surely the cavity identities and the bounds
\eqref{eq:site-field-bound} then hold simultaneously at all times and
all their indices.

With the response coefficients $c_{vu}$ from
\eqref{eq:covariance} and $A_v=\sum_u c_{vu}^2$ from
\eqref{eq:row-energy}, the continuous site field satisfies the equation
below.  In every integral, $A_v$ is assigned value zero wherever this
sum is infinite.
\begin{equation}
 \label{eq:field-sde}
 h_v(t)=\sum_{u\in\mathcal V}\int_0^t c_{vu}(s)\dd W_u(s)
            +\int_0^t m_v(s)A_v(s)\dd s.
\end{equation}
Each Brownian sum in this proposition converges, along any exhaustion
by finite vertex sets, in the norm
\[
 \|Z\|_{\mathcal M_T^2}
       =\left(\E\sup_{0\le t\le T}|Z(t)|^2\right)^{1/2}
       \qquad(T<\infty).
\]
The drift in \eqref{eq:field-sde} is absolutely integrable on each such
horizon almost surely and in expectation.
\end{proposition}

The raw and augmented filtrations both occur in the proof.  We record
the elementary passage between them once.  If $(\mathcal G_t^0)$ is a
raw filtration, write
$\mathcal G_t=\bigcap_{r>t}\overline{\mathcal G_r^0}$ for its usual
augmentation, where the bar includes the null sets of the ambient
probability space.

\begin{lemma}[Passage to the usual augmentation]
 \label{lem:continuous-augmentation}
If $Y$ is a continuous square-integrable martingale in
$(\mathcal G_t^0)$ and is continuous in $L^2$ at each time, then it is
also a martingale in $(\mathcal G_t)$.  More generally, if $Z$ is
integrable and a version $M(t)=\E[Z\mid\mathcal G_t^0]$ is continuous
in $L^1$, then $M(t)=\E[Z\mid\mathcal G_t]$ at every fixed time.
\end{lemma}

\begin{proof}
For $s<t$, choose $s_n\downarrow s$ with $s_n<t$.  The completed
sigma-fields $\overline{\mathcal G_{s_n}^0}$ decrease to
$\mathcal G_s$.  Reverse conditional expectation convergence gives
\[
 \E[Y(t)\mid\mathcal G_s]
   =\lim_n\E[Y(t)\mid\overline{\mathcal G_{s_n}^0}]
   =\lim_n Y(s_n)=Y(s)
\]
in $L^1$; the last limit even holds in $L^2$.  Completion does not
change conditional expectations up to null sets.  The same argument,
with $Z$ in place of $Y(t)$, proves the second assertion.  A continuous
bounded process is continuous in every finite $L^p$, by dominated
convergence, so the second assertion applies to the bounded posterior
versions constructed below.
\end{proof}

\subsection{Brownian innovations and finite posterior energy}

We first prove the asserted Brownian property.  For $0\le s<t$, the
future increments of the planted Brownian motions are independent of
all spins and of their Brownian pasts.  Conditioning on the smaller
raw observation history therefore gives
\[
 \E[X_v(t)-X_v(s)\mid\mathcal F_s^0]=(t-s)m_v(s).
\]
For every deterministic $r\ge s$, the tower property for the bounded
raw posteriors also gives
$\E[m_v(r)\mid\mathcal F_s^0]=m_v(s)$.  Conditional Fubini is
applicable because $|m_v|\le1$.  Subtracting the conditional means of
the two increments in \eqref{eq:innovations-definition} shows that
$W_v$ is a continuous square-integrable martingale in
$(\mathcal F_t^0)$.

This conclusion concerns the innovations, not the planted noises.
In the filtration containing all spins at time zero and all Brownian
histories, their relation is
\[
 W_v(t)=B_v(t)+\int_0^t(\sigma_v-m_v(s))\dd s.
\]
The added path has finite variation.  Hence the pathwise quadratic
covariations are
\[
 [W_v,W_u]_t=\mathbf 1_{\{v=u\}}t.
\]
The bounded drift and $L^2$ continuity of $B_v$ give $L^2$ continuity
of $W_v$.  Lemma~\ref{lem:continuous-augmentation} makes each
coordinate a martingale in $(\mathcal F_t)$.  L\'evy's characterization
applied to each finite vector, with the displayed covariations, proves
that vector is standard Brownian motion in $(\mathcal F_t)$ and hence
also in the smaller raw filtration.  Its future increments are
independent of the observation past.  By a monotone-class argument
over finite vertex sets and rational future times, the entire future
increment field
$(W_v(t+s)-W_v(t))_{v\in\mathcal V,s\ge0}$ is independent of
$\mathcal F_t$ for every fixed $t$.

Finite observation sets have their own innovations.  They will provide
an exact scalar energy identity, but we will never identify their
Brownian paths with the infinite innovations.

\begin{lemma}[Finite posterior energy]
 \label{lem:finite-energy}
Let $D$ be a finite connected subtree.  For $u\in D$ put
\[
 W_u^D(t)=X_u(t)-\int_0^t m_u^D(s)\dd s.
\]
The vector $W^D$ is standard Brownian motion in the usual augmentation
of $(\mathcal F_t^D)$, and, for $v\in D$,
\begin{equation}
 \label{eq:finite-energy}
 \begin{split}
 m_v^D(t)&=\sum_{u\in D}\int_0^t C_{vu}^D(s)\dd W_u^D(s),\\
 \E[m_v^D(t)^2]
   &=\E\int_0^t\sum_{u\in D}C_{vu}^D(s)^2\dd s .
 \end{split}
\end{equation}
\end{lemma}

\begin{proof}
The preceding conditional-increment and covariation calculation works
with the observations restricted to $D$, and proves the Brownian
assertion.  The finite posterior formula
\eqref{eq:finite-posterior} gives
$m_v^D(t)=M_v^D(X_D(t))$ for a smooth function $M_v^D$.
It is a bounded continuous martingale in the raw filtration, since it
is a version of $\E[\sigma_v\mid\mathcal F_t^D]$.  It is continuous in
$L^2$, so Lemma~\ref{lem:continuous-augmentation} also makes it a
martingale in the augmented finite observation filtration.

Apply finite-dimensional It\^o's formula to $M_v^D(X_D)$ using
$\dd X_u=m_u^D\dd t+\dd W_u^D$.  By
\eqref{eq:finite-posterior-derivative}, its continuous local martingale
part is $\sum_{u\in D}\int C_{vu}^D\dd W_u^D$.  The remaining part is
continuous and has finite variation locally.  Since $m_v^D$ is itself
a martingale, that part is also a local martingale and must vanish.
The covariance coefficients are bounded, so the finite Brownian sum
is square integrable on every finite horizon.  The free prior has mean
zero.  This proves the first identity in \eqref{eq:finite-energy}, and
orthogonality of the Brownian coordinates and It\^o's isometry prove
the second.
\end{proof}

\subsection{No covariance energy escapes to the boundary}

The finite energy identity does not yet imply its infinite analogue:
coordinatewise convergence could leave a positive amount of energy at
ever more distant vertices.  We control that possibility in two
intervals.  Away from time zero, each vertex's own Gaussian noise
attenuates its path factor, including at a finite-tree leaf.  Near time
zero, the whole-site information estimate bounds the entire finite
energy at once.

\begin{lemma}[Gaussian attenuation away from time zero]
 \label{lem:gaussian-path-tail}
For every $\delta>0$ there is $\kappa_\delta\in(0,1)$ such that, for
every finite connected subtree $D$, every $v,u\in D$, and every
$s\ge\delta$,
\begin{equation}
 \label{eq:gaussian-path-tail}
 \E[C_{vu}^D(s)^2]
     \le\theta^{2r}(1-\kappa_\delta)^r,
       \qquad r=\dist(v,u).
\end{equation}
The same inequality holds for the whole-tree covariance $C_{vu}(s)$.
\end{lemma}

\begin{proof}
Choose $\varepsilon_\delta>0$ with
$2\varepsilon_\delta/\sqrt{2\pi\delta}\le1/2$.  A Gaussian random
variable $Z$ of variance $s\ge\delta$, with any mean, has
$\P(|Z|\le\varepsilon_\delta)\le1/2$ by its density bound.  Since
\[
 a(z)^2\le a(z)
      \le 1-(1-\theta^2)\tanh^2z,
\]
it follows that
\[
 \E[a(Z)^2]\le1-\kappa_\delta,\qquad
 \kappa_\delta=\frac{1-\theta^2}{2}
                         \tanh^2\varepsilon_\delta\in(0,1).
\]
The estimate is uniform in the mean.

Let $v=v_0,\ldots,v_r=u$ be the path in $D$.  In the squared covariance
formula, discard the initial factor
$(1-(m_v^D)^2)^2\le1$.  At the fixed time $s$ the remaining factors
depend only on terminal observations.  Conditional on all spins and
all terminal Brownian coordinates other than $B_{v_1}(s)$, the cavity
recursion has the form
\[
 h_{v_1\to v_0}^D(s)=B_{v_1}(s)+z,
\]
where the shift $z$ is fixed: all branch fields in this recursion
exclude $v_1$.  The later fields
$h_{v_i\to v_{i-1}}^D$, $i>1$, also exclude $v_1$.  Conditional on the
spins, the terminal Brownian coordinates are independent Gaussians of
variance $s$.  Integrating $B_{v_1}(s)$ first therefore removes the
first squared factor at a cost at most $1-\kappa_\delta$, without
changing the later ones.  Integrate $B_{v_2}(s),\ldots,B_{v_r}(s)$ in
that order.  Iteration proves \eqref{eq:gaussian-path-tail}, even
conditional on the spins.  At a leaf the shift is simply its planted
spin drift; the argument still applies.  No infinite side branch is
used here.

For fixed $v,u,s$, exhaust the tree by finite connected subtrees
containing the path.  Bounded conditional-expectation convergence
from Lemma~\ref{lem:posterior} gives $C_{vu}^D(s)\to C_{vu}(s)$ in
$L^2$.  Passing to the limit proves the whole-tree inequality.
\end{proof}

Fix $v$ and let $D_R$ be its radius-$R$ ball.  Extend
$C_{vu}^{D_R}$ by zero for $u\notin D_R$.  For each fixed $u$ and
$T<\infty$,
\begin{equation}
 \label{eq:coordinate-covariance-convergence}
 C_{vu}^{D_R}\longrightarrow C_{vu}
 \quad\text{in }L^2([0,T]\times\Omega,\dd s\,\dd\P).
\end{equation}
Indeed, at each deterministic $s$, increasing-sigma-field convergence
applies to the bounded variables $\sigma_v$, $\sigma_u$, and
$\sigma_v\sigma_u$.  Their conditional means determine the
covariance.  The covariances are uniformly bounded, so dominated
convergence also integrates the resulting $L^2$ convergence in time.

When $b\theta^2<1$, the pointwise bound
$C_{vu}^{D_R}(s)^2\le\theta^{2\dist(v,u)}$ controls all spatial tails:
there are $n_r=db^{r-1}$ vertices at distance $r\ge1$, and
\[
 \sup_R\E\int_0^T\sum_{\dist(v,u)>L}C_{vu}^{D_R}(s)^2\dd s
 \le \frac db T\sum_{r>L}(b\theta^2)^r\longrightarrow0.
\]
The same deterministic estimate holds for the infinite coefficients.

At equality, fix $0<\delta\le\min\{T,1\}$.  The tower property gives
$m_v^{D_R}(\delta)=
\E[m_v(\delta)\mid\mathcal F_\delta^{D_R}]$.  Conditional Jensen,
the finite energy identity, and the whole-site estimate in
\eqref{eq:critical-information} yield
\begin{equation}
 \label{eq:early-energy}
 \E\int_0^\delta\sum_u C_{vu}^{D_R}(s)^2\dd s
   =\E[m_v^{D_R}(\delta)^2]
   \le\E[m_v(\delta)^2]\le C\sqrt\delta .
\end{equation}
On $[\delta,T]$, Lemma~\ref{lem:gaussian-path-tail} applies uniformly
in $R$, and now $n_r\theta^{2r}=d/b$.  Thus
\begin{equation}
 \label{eq:uniform-energy-tail}
 \sup_R\E\int_0^T
       \sum_{\dist(v,u)>L}C_{vu}^{D_R}(s)^2\dd s
 \le C\sqrt\delta+\frac db T
                    \sum_{r>L}(1-\kappa_\delta)^r .
\end{equation}
This also bounds the infinite tail.  To see this without assuming
its finiteness, first restrict that tail to a finite annulus,
pass $R\to\infty$ using
\eqref{eq:coordinate-covariance-convergence}, and then increase the
outer radius using monotone convergence.  First send $L\to\infty$
with $\delta$ fixed, so that the geometric term vanishes.  Then send
$\delta\downarrow0$.  The order matters: $\kappa_\delta$ may shrink
as $\delta$ tends to zero.  We have proved that the integrated spatial
tails of both the finite and the infinite rows vanish uniformly.

For a fixed $L$, the sum over $\dist(v,u)\le L$ is finite.
Equation~\eqref{eq:coordinate-covariance-convergence} passes its
squared energy to the limit in $L^1$.  The uniform tail estimates
then pass the complete energy to the limit.  At each fixed $t$ we
also have $m_v^{D_R}(t)\to m_v(t)$ in $L^2$.  Consequently
\eqref{eq:finite-energy} gives the exact infinite identity
\begin{equation}
 \label{eq:covariance-energy}
 \E[m_v(t)^2]
    =\E\int_0^t\sum_{u\in\mathcal V}C_{vu}(s)^2\dd s
        \qquad(t\ge0).
\end{equation}
At $t=0$ both sides are zero.  Only the scalar energies have been
passed to the limit; the distinct processes $W^{D_R}$ have played
no part in this limiting operation.

\subsection{Identifying the entire posterior martingale}

The energy identity gives the norm that a Brownian representation must
have.  A second identity identifies its inner product with the
posterior.  Together they determine the posterior itself and exclude
any additional martingale orthogonal to the innovations.

We first record the approximation that permits coefficients to be
used in this test.  It will also let the subsequent conditional-copy
construction transport individual stochastic integrals.

\begin{lemma}[Observation-adapted step approximation]
 \label{lem:observation-step-approximation}
For fixed $v,u$ and $T<\infty$, each of $C_{vu}$ and $c_{vu}$ is the
$L^2([0,T]\times\Omega,\dd s\,\dd\P)$ limit of bounded elementary
processes
\[
 G(s)=\sum_{j=0}^{N-1}\xi_j\mathbf1_{(t_j,t_{j+1}]}(s),
 \qquad 0=t_0<\cdots<t_N=T,
\]
where $\xi_j$ is a bounded Borel function of finitely many
observations at time $t_j$.  In particular each coefficient has a
predictable representative for stochastic integration.
\end{lemma}

\begin{proof}
For $R$ large enough to contain the path from $v$ to $u$, the finite
coefficient $C_{vu}^{D_R}$ is a bounded continuous function of the
terminal observations in $D_R$, and is therefore a continuous raw
adapted process.  It converges to $C_{vu}$ in the required space by
\eqref{eq:coordinate-covariance-convergence}.  For the response
coefficient use
\[
 c_{vu}^{D_R}=\theta^{\dist(v,u)}
       \prod_{i=1}^{\dist(v,u)}
                     a(h_{v_i\to v_{i-1}}^{D_R}).
\]
The component convergence in Lemma~\ref{lem:posterior} gives
fixed-time convergence to $c_{vu}$; the common bound
$|c_{vu}^{D_R}|\le\theta^{\dist(v,u)}\le1$ again gives convergence
in the stated $L^2$ space.  This finite coefficient is also a
continuous raw adapted function of finitely many terminal
observations.

For either finite coefficient, sample its values at the left
endpoints of a deterministic mesh.  As the mesh tends to zero,
continuity and the common bound imply convergence in
$L^2(\dd s\,\dd\P)$ by dominated convergence.  Choose a diagonal
sequence of radii and meshes.  These are the asserted elementary
processes.  A subsequence converges $\dd s\,\dd\P$-almost everywhere;
its pointwise limit where finite, with value zero elsewhere, is
predictable because each elementary process is predictable.  This
representative agrees with the original coefficient in the required
product-measure sense.
\end{proof}

Fix $v,u$ and $t<\infty$.  For any bounded elementary process of the
same form whose $\xi_j$ is merely
$\mathcal F_{t_j}^0$-measurable, we claim
\begin{equation}
 \label{eq:innovation-test}
 \E\left[m_v(t)\int_0^tG(s)\dd W_u(s)\right]
       =\E\int_0^tG(s)C_{vu}(s)\dd s.
\end{equation}
For an elementary process, the integral is the finite sum of
$\xi_j(W_u(t_{j+1})-W_u(t_j))$, hence is measurable in
$\mathcal F_t^0$.  In its expectation against $m_v(t)$ we may
therefore replace $m_v(t)$ by $\sigma_v$.  In the planted filtration,
\[
 \dd W_u(s)=\dd B_u(s)+(\sigma_u-m_u(s))\dd s.
\]
Here $\sigma_vG$ is a bounded predictable integrand: $\sigma_v$ is
known at time zero and every raw observation past is contained in the
planted past.  Thus
$\E[\sigma_v\int_0^tG\dd B_u]=0$.  For the finite-variation part,
condition at time $s$ on the raw observation history:
\[
 \E[\sigma_v(\sigma_u-m_u(s))\mid\mathcal F_s^0]
  =\E[\sigma_v\sigma_u\mid\mathcal F_s^0]-m_v(s)m_u(s)
  =C_{vu}(s).
\]
Boundedness permits Fubini, proving \eqref{eq:innovation-test}.
This argument has used only the fact that $W$ is a function of the
observation past, not any converse measurability.

Apply Lemma~\ref{lem:observation-step-approximation} to $C_{vu}$.
It\^o's isometry for $W_u$ and Cauchy--Schwarz on the two sides of
\eqref{eq:innovation-test} pass to the limit, so the identity also
holds with $G=C_{vu}$.  For a finite vertex set $F$ let
\[
 S_{v,F}(t)=\sum_{u\in F}\int_0^tC_{vu}(s)\dd W_u(s).
\]
The integrals for different $u$ are orthogonal, even though their
integrands may depend on every observation: their driving Brownian
coordinates have zero mutual quadratic covariation.  The isometry and
the testing identity give
\[
 \E[S_{v,F}(t)^2]
   =\sum_{u\in F}\E\int_0^t C_{vu}(s)^2\dd s
   =\E[m_v(t)S_{v,F}(t)] .
\]
Combining this with the exact energy
\eqref{eq:covariance-energy} yields the useful squared-error identity
\begin{equation}
 \label{eq:innovation-residual}
 \E[(m_v(t)-S_{v,F}(t))^2]
     =\sum_{u\notin F}\E\int_0^t C_{vu}(s)^2\dd s.
\end{equation}
The right side tends to zero as $F$ exhausts the vertices.  This is
precisely the step that eliminates a possible residual martingale.

There is also convergence of the processes, rather than just their
values at a fixed time.  For $F\subseteq F'$, Doob's maximal
$L^2$ inequality gives
\[
 \E\sup_{0\le s\le T}|S_{v,F'}(s)-S_{v,F}(s)|^2
 \le4\sum_{u\in F'\setminus F}
                    \E\int_0^T C_{vu}(s)^2\dd s .
\]
The total energy is finite by \eqref{eq:covariance-energy}.  The
partial sums therefore converge in $\mathcal M_T^2$ to a continuous
square-integrable martingale.  The limits agree on overlapping
horizons, and \eqref{eq:innovation-residual} identifies each
fixed-time value with the raw posterior $m_v(t)$.  Countability of
the vertices lets us choose these continuous versions simultaneously.
This proves \eqref{eq:magnetization-sde}.

\subsection{Continuous versions that retain component locality}
 \label{sec:continuous-causal-versions}

The Brownian representation has supplied a continuous version of each
site mean.  We now select Borel observation functions realizing this
version and continuous component fields.  Constructing component
fields from site fields could apparently use information from the
opposite side of an edge.  We first establish continuous component
modifications, then choose them by explicit past-observation formulas
that retain the required component restriction.

Write $M_v$ for the continuous mean version temporarily.  At every
rational time it agrees almost surely with the raw mean from
Lemma~\ref{lem:posterior}.  Intersect these events over all vertices
and rational times.  The raw field bound and continuity of $M_v$ and
$X_v$ imply, on this one event,
\[
 |M_v(t)|\le\tanh(|X_v(t)|+d\beta)<1\qquad(t\ge0).
\]
Indeed the inequality holds first at rational times and then passes
to every time by continuity.  On each compact interval,
\[
 \sup_{0\le t\le T}|M_v(t)|
 \le\tanh\left(\sup_{0\le t\le T}|X_v(t)|+d\beta\right)<1 .
\]
Thus $H_v=\operatorname{atanh}M_v$ is continuous.  The stronger
inequality $|H_v-X_v|\le d\beta$ also passes from rational times to
all times.

For adjacent vertices, consider for arbitrary real $H,K$ the equations
\[
 x+\psi(y)=H,\qquad y+\psi(x)=K.
\]
They have a unique solution.  In fact $x$ is the unique fixed point of
\[
 x\longmapsto H-\psi(K-\psi(x)),
\]
a contraction of the complete space $\mathbb R$ with constant at
most $\theta^2<1$, and $y=K-\psi(x)$.  Comparing two solutions gives
\[
 |x-\widetilde x|
 \le\frac{|H-\widetilde H|+\theta|K-\widetilde K|}{1-\theta^2},
\]
and the analogous bound for $y$.  The solution therefore depends
continuously on $(H,K)$.

Apply this construction to $(H_v(t),H_w(t))$ and denote the resulting
continuous processes by $(H^c_{v\to w},H^c_{w\to v})$.
At every fixed time the raw cavity identities identify this solution
almost surely with the raw component fields, by uniqueness.
After intersecting over rational times and oriented edges, continuity
therefore extends both cavity identities and the component bound to
all times.  This proves existence of continuous component
modifications.  The construction by adjacent site fields is used only
for that existence assertion.

We now specify the observation functions actually used in the rest of
the proof.  Let $m_v^0(t,x)$ and $h^0_{v\to w}(t,x)$ be the original
raw functions obtained from the centered finite-ball limits in
Lemma~\ref{lem:posterior}.  For
$q_n(t)=2^{-n}\lfloor2^nt\rfloor$, set
\[
 \widehat m_v(t,x)=
 \begin{cases}
  \displaystyle\lim_{n\to\infty}m_v^0(q_n(t),x),
       &\text{if the limit exists in }(-1,1),\\
  0,&\text{otherwise},
 \end{cases}
\]
and set $\widehat h_v=\operatorname{atanh}\widehat m_v$.  For every
oriented edge, set
\[
 \widehat h_{v\to w}(t,x)=
 \begin{cases}
  \displaystyle\lim_{n\to\infty}h^0_{v\to w}(q_n(t),x),
       &\text{if the limit exists in }\mathbb R,\\
  0,&\text{otherwise}.
 \end{cases}
\]
For each $n$ the sampled process is constant on the half-open dyadic
time intervals, and its value there is a Borel function of the
appropriate observation history at the left endpoint.  It is
therefore progressively measurable.  Real convergence is a countable
Cauchy condition.  For the site selector, membership of the resulting
limit in $(-1,1)$ is an additional Borel condition.  Taking these
limits with the displayed defaults therefore preserves Borel and
progressive measurability.
For the component expression both the entire sequence and its default
decision use only observations in $T_{v\to w}$ through time $t$.
In particular, two inputs agreeing on that component history have
the same sequence and the same chosen value.  This is pointwise
component locality, including on inputs where the limit fails.

At every dyadic time the raw site mean agrees almost surely with $M_v$,
and the raw component field agrees almost surely with $H^c_{v\to w}$.
Intersecting over the countable dyadic times and indices gives one
event of probability one.  On that event, for every real $t$ the
sampled values above are values of the respective continuous
modifications at $q_n(t)$.  Since $q_n(t)\to t$, continuity makes
their limits agree simultaneously for every $t$.  Thus the chosen
functions are indistinguishable from the continuous modifications.
They retain all the all-time cavity identities and bounds just
proved.  The same dyadic rule and the same defaults at every index
also preserve the graph-relabeling covariance of the raw functions.

We henceforth use these chosen versions and omit the hats.  Define
$c_{vu}$ again by its finite path product and
$C_{vu}=(1-m_v^2)c_{vu}$.  These are continuous in time almost surely,
and they agree at each fixed time with their former raw versions.
Consequently every changed coefficient agrees with its former version
outside a $\dd t\,\dd\P$-null set.  Likewise, Fubini and boundedness
give equality of the old and new mean integrals on each integer
horizon, and hence, after a countable intersection,
\[
 \int_0^t m_v^0(s,X)\dd s=\int_0^t m_v(s,X)\dd s
 \quad\text{simultaneously for all }t\ge0
 \quad\text{almost surely}.
\]
Thus the innovation paths have not changed.

These versions are also the posteriors for the usual filtrations.
For a fixed time $t$, apply the second part of
Lemma~\ref{lem:continuous-augmentation} to the bounded continuous
version of $\E[\sigma_v\mid\mathcal F_t^0]$.  It gives
$m_v(t)=\E[\sigma_v\mid\mathcal F_t]$.  The process
$C_{vu}+m_vm_u$ is a bounded continuous version of the raw conditional
mean of $\sigma_v\sigma_u$; applying the same argument identifies
it with that conditional mean in $\mathcal F_t$.  This identifies
$C_{vu}$ as the augmented conditional covariance.  Applied instead
to $\tanh h_{v\to w}$ and the raw component filtration, the argument
gives the corresponding assertion for its usual augmentation.
These statements include time zero.  The one-spin site and component
means start at zero.  The two-spin posterior starts at its prior value
$\E[\sigma_v\sigma_u]=\theta^{\dist(v,u)}$, including value one when
$u=v$.

\subsection{The field equation}

It remains to pass from the posterior mean to its field.  This requires
the response row $A_v=\sum_u c_{vu}^2$, rather than merely the
covariance row: the transformation divides by $1-m_v^2$.  The
estimate \eqref{eq:row-energy}, proved before this section, gives
\[
 \E\int_0^T A_v(s)\dd s<\infty
 \qquad(T<\infty).
\]
Its strict case follows from a deterministic geometric sum.  At
criticality it follows from the intact infinite-side-branch estimate
$\E A_v(s)\le C/\min\{\sqrt s,1\}$ for $s>0$.  This is a different
use of attenuation from the finite-boundary Gaussian estimate above.
The inputs here are the second-moment information and forward
infinite-branch attenuation already proved.

For clarity, define the version used in the drift directly from the
chosen observation functions:
\[
 A_v(s,x)=
 \begin{cases}
  \displaystyle\sum_{u\in\mathcal V}c_{vu}(s,x)^2,
      &\text{if this sum is finite},\\
  0,&\text{if this sum is infinite}.
 \end{cases}
\]
The extended nonnegative sum is a countable sum of Borel progressive
functions, so this finite-or-zero version is again a causal Borel
progressive observation function.  Permuting the summation index and
using the same default preserve graph relabeling.  The integrated row
estimate shows that this version agrees with the extended sum outside
a $\dd s\,\dd\P$-null set.  The predictable representatives used for
Brownian integration agree with the chosen path functions outside one
such null set for all $u$, by countability.  Their squared sum
therefore agrees with this $A_v$ there.  At criticality the unmodified
row is infinite at time zero, as noted after \eqref{eq:row-energy};
the displayed convention supplies the value used in the drift.

The Brownian representation and the covariance formula now imply, as
an equality of densities up to $\dd s\,\dd\P$-null sets,
\[
 \dd\langle m_v\rangle_s
    =\sum_u C_{vu}(s)^2\dd s
    =(1-m_v(s)^2)^2 A_v(s)\dd s.
\]
One may obtain the first equality from finite Brownian sums: their
brackets converge in $L^1$ total variation on every finite horizon,
by \eqref{eq:covariance-energy}, while their martingales converge in
$\mathcal M_T^2$.  The defining bracket identity then passes to the
limit.

For $n\ge2$ let
\[
 \tau_n=n\wedge\inf\{t\ge0:|m_v(t)|\ge1-1/n\}.
\]
On the stopped interval the derivatives
\[
 (\operatorname{atanh})'(x)=\frac1{1-x^2},\qquad
 (\operatorname{atanh})''(x)=\frac{2x}{(1-x^2)^2}
\]
are bounded.  It\^o's formula for the stopped continuous martingale
therefore gives
\[
 h_v(t\wedge\tau_n)
   =\sum_u\int_0^{t\wedge\tau_n}
             \frac{C_{vu}(s)}{1-m_v(s)^2}\dd W_u(s)
       +\int_0^{t\wedge\tau_n}m_v(s)A_v(s)\dd s.
\]
The Brownian coefficient is exactly $c_{vu}$.  The equality of this
countable Brownian sum with the transformed mean integral follows
first for finite sums and then by the isometry; on the stopped
interval the derivative is bounded.

The global Brownian sum with coefficients $c_{vu}$ exists in
$\mathcal M_T^2$.  Indeed, for finite $F\subseteq F'$, Doob's inequality
gives
\[
 \E\sup_{0\le t\le T}
 \left|\sum_{u\in F'\setminus F}\int_0^t c_{vu}(s)\dd W_u(s)\right|^2
 \le4\sum_{u\in F'\setminus F}\E\int_0^T c_{vu}(s)^2\dd s.
\]
The right side tends to zero along an exhaustion by the integrated
row estimate.  Also
$\E\int_0^T|m_v|A_v\dd s\le\E\int_0^T A_v\dd s<\infty$.
Finally $\tau_n\uparrow\infty$ almost surely: the continuous mean stays
strictly inside $(-1,1)$ on every compact interval.  The stopped
identities hold simultaneously in time by continuity.  On any fixed
compact interval, almost every path therefore agrees with one of
these identities for all sufficiently large $n$.  Removing the
stopping yields \eqref{eq:field-sde}.

We now have the full innovation-driven field equation in the observation
filtration, with both Brownian and drift terms integrable.  The next
comparison must prove that the same innovations remain Brownian when
two conditional observation histories are put together; that
compatibility is the next step of the argument.
